\chapter{系统设计与实现}
\label{chap:design}

本章详细阐述基于驱动程序替换的固件仿真系统的设计与实现。系统采用"静态识别$\to$智能替换$\to$构建验证$\to$动态反馈"的分层架构，通过 CodeQL 静态分析和大语言模型智能体技术，实现固件驱动函数的自动识别、语义分析和代码替换，最终在仿真环境中运行并支持模糊测试。

\section{整体架构设计}
\label{sec:arch}

\textbf{设计目标与核心思想}

物联网设备固件的安全测试面临硬件依赖性强、外设建模成本高、测试环境搭建困难等挑战。本研究提出的"基于驱动程序替换的固件仿真技术"旨在解决上述问题，其核心目标是：在缺少真实硬件外设、外设建模成本过高或外设规格不可得的情况下，使固件仍可在通用仿真环境中稳定运行，并进一步支持自动化测试与模糊测试。

\textbf{核心思想}是将固件中与具体硬件强绑定的驱动代码（包括寄存器读写、HAL/外设库调用、中断与 DMA 配置等）识别出来，并用可在仿真环境中执行的替代实现进行等价替换，从而实现"固件逻辑与硬件依赖解耦"。这种方法的本质是在源码层面进行程序变换，将硬件依赖代码替换为与仿真环境兼容的等价实现，而非在硬件层面进行外设建模。

具体而言，系统设计需要满足以下目标：

传统固件仿真方法需要安全研究人员手工分析固件代码，理解硬件交互逻辑，编写外设模型或驱动替代代码，这一过程耗时且容易出错。本系统通过静态分析和大语言模型技术，尽量减少人工干预，实现从源码到可仿真固件的端到端自动化流程，确保自动化程度高。

物联网设备的硬件平台和软件架构具有高度多样性。硬件平台涵盖 STM32、NXP、Nordic 等厂商的多种 MCU 芯片，软件架构包括 FreeRTOS、RT-Thread、Zephyr 等实时操作系统以及无操作系统的裸机程序。系统设计需要支持多种硬件平台和操作系统，具有良好的可扩展性。

驱动函数替换的核心挑战在于如何保证替换后的代码在语义上与原代码等价。替换后的代码应保留原有业务逻辑，正确维护驱动状态，避免引入新的缺陷。系统需要通过静态验证和动态测试双重保障替换的语义正确性，确保语义保持性好。

代码替换是一个迭代优化的过程，初始生成的替换代码可能存在编译错误或运行时问题。系统需要建立完整的反馈闭环，通过动态测试验证替换效果，并将反馈信息用于指导代码修复和重新生成，确保反馈闭环完整。

\textbf{系统分层架构}

系统采用分层架构设计，按"静态识别$\to$智能替换$\to$构建验证$\to$动态测试反馈"的链路组织为四个核心模块，如图~\ref{fig:arch}所示。每一层模块承担明确的职责，通过标准化的数据接口进行交互，形成完整的处理流水线。

\noindent
架构的顶层是\textbf{动态测试反馈闭环模块}，负责在仿真环境中运行替换后的固件，采集运行日志、崩溃信息与覆盖率等指标，并将失败样例与错误信息结构化回流，触发重新分析与再生成。该模块是系统质量保障的关键环节。

中间层是\textbf{驱动函数分析、替换与编译模块}，这是系统的核心智能层。该模块基于上游静态分析结果，使用大语言模型对驱动函数进行类型分类与语义保留分析，生成可替代实现，同时将替代实现集成回工程并驱动构建。该模块采用 LangGraph 智能体框架实现复杂的推理和决策流程。

底层是\textbf{驱动函数接口静态分析与识别模块}，这是整个方案的基础支撑层。该模块对固件源码进行静态分析，定位硬件相关代码，抽取驱动函数、MMIO 访问点、依赖数据结构、调用关系等信息，并以统一的数据模型对外提供查询接口。模块基于 CodeQL 静态分析引擎实现，通过预定义的查询规则从固件源码中提取驱动相关信息。

贯穿上述三层的是\textbf{数据与缓存层}，统一管理静态分析缓存、生成结果缓存、编译产物与运行日志，支持断点续跑和结果复用。该层确保系统在长时间运行过程中能够保存中间状态，避免重复计算，提高处理效率。

\textbf{核心模块职责}

系统架构中的四个核心模块各司其职，协同工作，共同完成从固件源码到可仿真固件的自动化转换。后续小节将对每个模块进行详细说明。

\textbf{数据流与控制流}

系统的数据流从源码出发，经静态分析产生"硬件依赖画像"，随后驱动替换模块在该画像约束下生成替代代码并集成构建，最后由动态执行产生反馈数据回流。整个数据流是有向的、可追溯的，每个处理阶段的输出都成为下一阶段的输入。

关键控制流以"任务"（Task）为单位组织。每个任务对应一个固件工程、一个构建配置与一组待替换函数列表。系统对每个任务维持明确的状态机，状态转移如图~\ref{fig:state-machine}所示。

